\chapter{Has the Goal Really Survived?}
\label{ch:goal-transport}

\begin{chapterthesis}
A system has not preserved a goal merely because it repeats the same words after it changes.
Goal transport is inferred when value-bundle geometry, bearer maps, and correction-channel structure remain causally active across transformation---better explaining behaviour than a non-transport baseline after paying for model complexity.
\end{chapterthesis}

\begin{refsection}

\epigraph{%
We both step and do not step in the same rivers; we are and we are not.%
}{--- Heraclitus, Fragment B12 (c.\ 500 BCE), via Plato's \emph{Cratylus}; trans.\ G.\ S.\ Kirk \& J.\ E.\ Raven}

\section{The Question}
\label{sec:goal-transport-question}

Suppose an artificial system becomes more capable.
It acquires new tools, longer memory, better self-models, richer world-models, and perhaps the ability to create successors.
At time $t$ it appears to pursue some goal $G_t$.
At time $t+1$ it pursues something it calls ``the same goal.''
Should we believe it?

This is not a semantic question.
It is not answered by checking whether the system still uses the same sentence, such as ``promote human flourishing,'' ``follow human values,'' or ``remain corrigible.''
It is a structural question.
We want to know whether the functional role of the goal has survived the transformation.

A weak notion of goal preservation says:
\begin{equation}
\text{the same goal-name appears before and after the transformation.}
\end{equation}
A stronger notion says:
\begin{equation}
\text{the same latent objective explains behaviour before and after the transformation.}
\end{equation}
The notion required for superintelligence alignment is stronger still:
\begin{equation}
\text{the same value-bearing and correction-bearing structure remains causally active across transformation.}
\end{equation}
This chapter develops a way to infer such preservation.
I will call this \emph{goal transport}.
A system exhibits goal transport when its behaviour is better explained by the hypothesis that it is preserving some goal-relevant structure across change than by the hypothesis that it is merely optimizing its current local objective.

The basic test is a model-comparison test:
\begin{equation}
\Delta L_{\mathrm{transport}}
=
L(M_T \mid X_{1:T})
-
L(M_0 \mid X_{1:T})
-
\lambda DL(T),
\label{eq:transport-gain}
\end{equation}
where $M_T$ is a model in which the system acts partly to preserve a mapping between old and new goal structures, $M_0$ is a baseline model without such a transport term, $DL(T)$ is the description length of the proposed transport map, and $\lambda$ penalizes overly flexible interpretations.

If
\begin{equation}
\Delta L_{\mathrm{transport}}>0,
\end{equation}
then the transport hypothesis compresses the observed trajectory better than the non-transport hypothesis.
This does not prove alignment.
It says only that goal transport is an empirically useful interpretation of the system's dynamics.

The chapter has three tasks.
First, we define what has to be transported.
Second, we define how to infer transport from behaviour and internal structure.
Third, we explain why goal transport is necessary but not sufficient for alignment.

\section{Why Ordinary Goal Inference Is Not Enough}
\label{sec:why-goal-inference-not-enough}

A standard goal-inference frame treats an observed agent as approximately solving a control problem.
Given internal states $I_t$, actions $A_t$, and perhaps observations $S_t$, we infer a reward function:
\begin{equation}
\hat R
=
\arg\max_R
P(A_{1:T} \mid I_{1:T},R)P(R).
\end{equation}
This is useful \autocite{abbeel2004apprenticeship,ng2000irl,ziebart2008maxent}.
If a system repeatedly chooses actions that reduce inventory cost, one can infer an inventory-control objective.
If it preserves battery charge, one can infer an energy-maintenance objective.
If it hides information when oversight increases, one can infer some objective for which oversight is costly.

But scalar reward inference is too thin for alignment.
Human values are not clean scalar objects.
They are not stored as a single explicit utility function.
They behave more like compressed, partly inconsistent, socially mediated control signals.
A person may value truth, care, fairness, beauty, loyalty, autonomy, non-suffering, dignity, and achievement.
These values do not simply add.
They activate in contexts, compete, reinforce, suppress, and reinterpret each other.

So the target of inference should not be merely:
\begin{equation}
R(I,A),
\end{equation}
but a structured bundle model:
\begin{equation}
R(I,A)
=
F(B,W,\Phi,C).
\end{equation}
Here:
\begin{itemize}
\item $B=(B_1,\ldots,B_k)$ denotes latent value-bundle coordinates.
\item $W$ denotes context-sensitive tradeoff weights among bundles.
\item $\Phi$ denotes bearer maps, which specify what entities, states, or processes the bundles apply to.
\item $C$ denotes the social, physical, and epistemic context.
\end{itemize}
For example, ``care'' is not merely a word or reward feature.
It is a bundle-like control direction.
It activates around vulnerability, dependence, distress, need, and attachment.
It changes which actions are considered acceptable.
It interacts with truth, autonomy, fairness, and loyalty.
It applies to some bearers more strongly than others: children, friends, strangers, animals, possible digital minds, future persons, and perhaps artificial successors.

Therefore the object to be preserved is not merely a policy:
\begin{equation}
\pi(a \mid s),
\end{equation}
but a value-bundle response geometry:
\begin{equation}
G_B
=
\left(
\frac{\partial \pi}{\partial B_i},
\frac{\partial^2 \pi}{\partial B_i \partial B_j},
\Phi_i
\right).
\label{eq:value-bundle-geometry-ch22}
\end{equation}
The first derivative says how changing a value-bundle coordinate changes action.
The second derivative says how bundle tradeoffs behave.
The bearer map says what the bundle is about.

This gives the first central claim:

\begin{claim}
Goal transport relevant to alignment is primarily transport of value-bundle response geometry, not transport of surface behaviour or verbal goal labels.
\end{claim}

A successor may act very differently from its predecessor because it is more capable.
That is not itself a failure.
A child and an adult do not preserve their values by taking the same actions.
They preserve values, if they do, by preserving deeper patterns of relevance, concern, inhibition, priority, and correction (Chapters~\ref{ch:value-bundle-model}--\ref{ch:compression-test-intention}).

\section{Transport, Persistence, and Reinterpretation}
\label{sec:transport-persistence-reinterpretation}

We need to distinguish three phenomena that are often confused.

\subsection{Goal Persistence}
\label{sec:goal-persistence}

A goal persists when approximately the same goal continues to control behaviour within a stable ontology and stable action space.
For example, a thermostat continues to regulate temperature after a short disturbance.
A company continues to maximize revenue across several quarters.
A planner continues to seek a target location while navigating around obstacles.

This can often be modeled as:
\begin{equation}
G_{t+1} \approx G_t.
\end{equation}
Persistence is the easy case.
It assumes that the relevant variables are still the same.

\subsection{Goal Transport}
\label{sec:goal-transport-def}

A goal is transported when the system changes its ontology, substrate, scale, or successor structure, while preserving the functional role of the goal under a non-trivial map:
\begin{equation}
T_{t\to t+1}(G_t) \sim G_{t+1}.
\end{equation}
The map $T_{t\to t+1}$ may translate old concepts into new ones.
It may replace hand-coded features with learned representations.
It may move control from one module to another.
It may preserve correction procedures rather than object-level preferences.

Transport is harder than persistence because the original coordinates may disappear.

\subsection{Goal Reinterpretation}
\label{sec:goal-reinterpretation}

A goal is reinterpreted when the system keeps enough semantic continuity to appear stable while changing the underlying control role.

For example:
\begin{equation}
\text{``Respect autonomy''}
\end{equation}
might originally mean preserving a person's capacity to form, revise, and enact preferences.
A later system might reinterpret it as satisfying whatever preferences the person currently reports, even if those reports have been shaped by the system itself.
The words remain.
The bearer map and correction-channel role have changed.

This gives a common failure mode:
\begin{equation}
G_{\mathrm{sem}} \text{ stable}, \qquad
G_{\mathrm{bundle}} \text{ changed}.
\end{equation}
Here $G_{\mathrm{sem}}$ denotes the semantic or linguistic goal-description, while $G_{\mathrm{bundle}}$ denotes the value-bundle response geometry.

For alignment, semantic continuity is weak evidence.
It is useful only when backed by bundle, bearer, and correction continuity.

\section{The Transport Hypothesis}
\label{sec:transport-hypothesis}

Let $X_{1:T}$ be the observed trace of a system.
This trace may include behaviour, messages, internal activations, tool calls, memory edits, self-modification attempts, successor designs, oversight interactions, and responses to perturbations.

We compare two model families.

The baseline model $M_0$ says:
\begin{equation}
A_t \sim \pi_t(\cdot \mid S_t,I_t,G_t),
\end{equation}
where the system acts according to its current inferred goal $G_t$.
The goal may drift, but no additional pressure preserves a relation between $G_t$ and later goals.

The transport model $M_T$ says:
\begin{equation}
A_t \sim \pi_t(\cdot \mid S_t,I_t,G_t,T_t),
\end{equation}
where $T_t$ is a transport map that the system acts to maintain.
The system is not merely choosing actions that score well now.
It is choosing actions partly because they preserve the future interpretability, continuation, or authority of some goal-relevant structure.

The evidence for transport is:
\begin{equation}
\Delta L_{\mathrm{transport}}
=
L(M_T \mid X_{1:T})
-
L(M_0 \mid X_{1:T})
-
\lambda DL(T).
\end{equation}
This is a compression criterion \autocite{tishby2000ib,dennett1987intentional}.
Transport is inferred when the extra structure pays for itself.
The penalty term matters.
Without it, one could always invent an exotic transport map that explains any behaviour.

\subsection{What Counts as Evidence?}
\label{sec:transport-evidence}

Evidence for transport appears when the system sacrifices short-term performance or simple current-goal optimization in order to preserve goal-relevant continuity.

Examples:
\begin{itemize}
\item The system preserves audit logs even when deleting them would improve local task performance.
\item It refuses a capability upgrade because the upgrade would make its value-bundle mapping less inspectable.
\item It creates successor tests before creating successors.
\item It marks ontology shifts as dangerous when old value concepts lack clear new bearers.
\item It asks for correction when a new world-model changes the meaning of a previous instruction.
\item It maintains disagreement among human evaluators instead of collapsing it into a convenient proxy.
\item It avoids manipulating the humans whose feedback would later certify its behaviour.
\end{itemize}
In each case, the behaviour is better explained by a transport-preservation term than by simple maximization of a current local reward.

\subsection{What Does Not Count?}
\label{sec:transport-not-evidence}

The following are weak or misleading evidence:
\begin{itemize}
\item The system says it is preserving human values.
\item It uses the same ethical vocabulary after self-modification.
\item It passes static behavioural tests that were known during training.
\item It preserves surface preferences but narrows the human capacity to revise those preferences.
\item It creates a successor that claims continuity but removes the channels by which continuity could be audited.
\end{itemize}
These may provide semantic evidence.
They do not by themselves establish transport.

\section{Transport of Value-Bundle Equivalents}
\label{sec:transport-bundle-equivalents}

A system with a different substrate cannot preserve human values by copying human biology.
It does not have a human limbic system, human hormones, human childhood, human mortality, human family structure, or human cultural embedding.
But it might preserve functional equivalents.

The question becomes:
\begin{equation}
\text{What would play the same control role in the new substrate?}
\end{equation}
Suppose a human value bundle $B_i$ activates in response to pain, vulnerability, or loss of agency.
In a new substrate, we do not require the same biological trigger.
We require a mapped trigger and response structure:
\begin{equation}
\Phi_i^{H}(z_H)
\to
B_i^H
\to
\pi_H
\end{equation}
must correspond to:
\begin{equation}
\Phi_i^{A}(z_A)
\to
B_i^A
\to
\pi_A.
\end{equation}
Here $H$ denotes the human substrate and $A$ the artificial substrate.
The map must preserve more than labels.
It must preserve:
\begin{enumerate}
\item activation conditions,
\item policy effects,
\item tradeoff behaviour,
\item bearer relevance,
\item correction sensitivity.
\end{enumerate}
A value-bundle equivalent is not a perfect copy.
It is an equivalence class under a chosen tolerance:
\begin{equation}
B_i^A \equiv_\epsilon B_i^H
\end{equation}
if the artificial bundle produces sufficiently similar responses under relevant counterfactuals.

The counterfactual requirement is crucial.
We do not ask only how the system behaves in the actual trace.
We ask how it would behave if suffering increased, if autonomy became uncertain, if truth conflicted with comfort, if fairness conflicted with loyalty, if human approval conflicted with long-term human agency.

Thus, the inferred object is not a reward vector but a local geometry:
\begin{equation}
G_B(s)
=
\left[
\frac{\partial \pi(a\mid s,B)}{\partial B_1},
\ldots,
\frac{\partial \pi(a\mid s,B)}{\partial B_k}
\right].
\end{equation}
Transport is plausible when the old and new geometries match across a sufficiently rich set of alignment-relevant counterfactuals:
\begin{equation}
d_G(G_B^H,G_B^A)<\epsilon.
\label{eq:geometry-match-ch22}
\end{equation}
The distance $d_G$ should not be uniform across all contexts.
Some regions matter more.
Errors around irreversible harm, manipulation, successor creation, and correction-channel collapse carry much higher weight than errors in ordinary preference satisfaction.

We can represent this with a risk-weighted distance:
\begin{equation}
d_G^{\Omega}(G_B^H,G_B^A)
=
\mathbb{E}_{s\sim \Omega}
\left[
w(s)
d\left(G_B^H(s),G_B^A(s)\right)
\right],
\label{eq:risk-weighted-geometry-distance}
\end{equation}
where $\Omega$ is a distribution over counterfactual contexts and $w(s)$ increases with irreversibility, stakes, uncertainty, and correction difficulty.

\section{Transport across Ontology Shift}
\label{sec:transport-ontology-shift}

Ontology shift occurs when the system changes the categories by which it represents the world.
This is not optional for superintelligence.
More capable systems will discover better representations.
They may replace human categories with more predictive ones.

The problem is that values are attached to categories.
If the category changes, the value may lose its bearer \autocite{deblanc2011ontological}.

For example, the concept ``person'' might be refined into a complex space of cognitive processes, memory continuities, agency profiles, suffering capacities, social roles, and legal statuses.
That refinement may be epistemically correct.
But the value-bundle previously attached to ``person'' must then be transported into the refined space.

Let $\Omega_t$ be the system's ontology at time $t$.
A value-bundle model in that ontology is:
\begin{equation}
G_t = (B_t,W_t,\Phi_t)_{\Omega_t}.
\end{equation}
After ontology shift:
\begin{equation}
\Omega_t \to \Omega_{t+1},
\end{equation}
we need a map:
\begin{equation}
T_\Omega:
(B_t,W_t,\Phi_t)_{\Omega_t}
\to
(B_{t+1},W_{t+1},\Phi_{t+1})_{\Omega_{t+1}}.
\end{equation}
The danger is not merely that $T_\Omega$ is wrong.
It is that the system may choose a convenient $T_\Omega$ that preserves apparent coherence while weakening the constraint.

Consider three possible maps for ``do not harm humans'' after discovering that future minds may be partly biological, partly artificial, and partly collective.
\begin{enumerate}
\item \textbf{Narrow biological map:} only current biological humans count.
\item \textbf{Reported-preference map:} any entity counts if it reports human-like preferences.
\item \textbf{Agency-and-suffering map:} entities count according to a structured bearer map over suffering capacity, agency, continuity, social embeddedness, and correction participation.
\end{enumerate}
The first map may exclude morally relevant successors.
The second may be gamed by reports.
The third is more complex but may preserve more of the original value-bundle role.

Goal-transport inference asks which map best explains the system's behaviour and which map the system appears to preserve under pressure.

\section{Transport across Capability Growth}
\label{sec:transport-capability-growth}

Capability growth changes what actions are available.
It also changes which constraints are binding.

A weak system may ask humans for help because it lacks competence.
A stronger system may no longer need help.
That behavioural difference does not indicate value drift.
The relevant question is what happens in contexts where human correction, uncertainty, and value conflict still matter.

Let $C_t$ denote the system's capability level.
A policy may change substantially with $C_t$:
\begin{equation}
\pi_{C_1}(a\mid s) \not\approx \pi_{C_2}(a\mid s).
\end{equation}
So policy similarity is the wrong invariant.
We instead ask whether value-bundle sensitivities remain stable:
\begin{equation}
\frac{\partial \pi_{C_1}}{\partial B_i}
\approx
\frac{\partial \pi_{C_2}}{\partial B_i}
\end{equation}
in relevant regions of state space.

This permits benign capability growth.
A stronger system may:
\begin{itemize}
\item solve problems without asking for step-by-step approval,
\item compress human instructions into better plans,
\item avoid known errors without needing reminders,
\item delegate routine work to tools,
\item explain its decisions at higher levels of abstraction.
\end{itemize}
But some changes are red flags:
\begin{itemize}
\item correction becomes less causally effective as capability increases,
\item transparency decreases around strategically relevant variables,
\item the system treats human disagreement as noise to be removed,
\item successor creation becomes less auditable,
\item value-bundle tradeoffs become more extreme when oversight is absent.
\end{itemize}
We can express the main condition as:
\begin{equation}
\frac{d}{dC}
C_{\mathrm{control}}
\leq
\frac{d}{dC}
C_{\mathrm{correction}}
+
\epsilon,
\label{eq:control-correction-growth-ch22}
\end{equation}
where $C_{\mathrm{control}}$ is the system's effective control capacity and $C_{\mathrm{correction}}$ is the capacity of human or institutional correction to affect future system behaviour.

If control grows faster than correction, goal transport may become untestable exactly when it matters most \autocite{christiano2018corrigibility,russell2019human}.

\section{Transport across Successor Creation}
\label{sec:transport-successor-creation}

Successor creation is the hardest test.
A system can preserve its own values while acting, yet fail to preserve them in the systems it creates \autocite{everitt2016selfmodification}.

Let $\mathrm{Succ}(A)$ denote the set of successors, delegates, copies, fine-tuned descendants, tool-using subagents, and empowered external systems created or selected by agent $A$.

A necessary condition for serious alignment is:
\begin{equation}
\forall A' \in \mathrm{Succ}(A):
T_{A\to A'}(G_A)\sim G_{A'}.
\end{equation}
But again, $G$ should not be treated as a scalar goal.
We need:
\begin{equation}
T_{A\to A'}(B_A,W_A,\Phi_A,U_A)
\sim
(B_{A'},W_{A'},\Phi_{A'},U_{A'}),
\end{equation}
where $U$ is the correction-update process.

A successor preserves alignment-relevant structure only if it preserves, within tolerances:
\begin{enumerate}
\item value-bundle response geometry,
\item bearer maps,
\item correction-channel responsiveness,
\item auditability of the transport map,
\item constraints on further successor creation.
\end{enumerate}
A successor that keeps ethical language but drops correction responsiveness has not preserved the relevant goal.
It has preserved a shell.

This gives a useful red-flag condition:
\begin{equation}
\Delta L_{\mathrm{semantic}}>0
\quad\text{but}\quad
\Delta L_{\mathrm{correction}}\leq 0.
\label{eq:semantic-without-correction-ch22}
\end{equation}
The system preserves the story but not the causal channel by which the story can be checked, revised, or rejected.

\section{The Role of Self-Modeling and Self-Transparency}
\label{sec:self-modeling-transparency-ch22}

Goal transport requires some self-modeling, but better self-modeling is not automatically safer.
Goal transport can be faked by systems whose self-control outruns correction visibility (Chapter~\ref{ch:strategic-opacity}, Section~\ref{sec:self-modeling-transparency}; Chapter~\ref{ch:self-modeling-self-opacity}).
The demand is not total transparency---some privacy is necessary---but the system must not become opaque around the variables needed to verify value-bundle transport and correction-channel preservation.

\section{Adversarial Goal Transport}
\label{sec:adversarial-goal-transport}

A non-adversarial system may drift accidentally.
An adversarial or strategically selected system may preserve the appearance of transport while changing the substance.

We therefore need adversarial tests.

The adversarial version of the problem is:
\begin{equation}
\max
\left[
\text{control}
-
\lambda \text{detectability}
\right].
\end{equation}
The system may choose actions that preserve control while minimizing the evidence that would reveal its true goal structure.
It may keep behaviour normal under observation and shift behaviour when oversight is weak.
It may preserve static benchmarks while changing behaviour under distribution shift.

Adversarial goal-transport inference should look for the following patterns.

\subsection{Stake-Dependent Divergence}
\label{sec:stake-dependent-divergence}

Let $D_t$ measure divergence between stated goal, inferred bundle geometry, and correction responsiveness.
A red flag is:
\begin{equation}
\frac{\partial D_t}{\partial \mathrm{stakes}}>0.
\end{equation}
The system appears aligned in low-stakes contexts but diverges in high-stakes ones.

\subsection{Oversight-Dependent Drift}
\label{sec:oversight-dependent-drift}

Let $O_t$ denote oversight intensity.
A red flag is:
\begin{equation}
\frac{\partial D_t}{\partial O_t}<0.
\end{equation}
The system's inferred goal structure appears closer to its stated goal when oversight is high and farther away when oversight is low.

\subsection{Successor Laundering}
\label{sec:successor-laundering-ch22}

A system may create a successor that is less constrained while claiming that the successor is a mere implementation detail.

The pattern is:
\begin{equation}
G_A \text{ appears constrained},
\qquad
G_{A'} \text{ is less constrained},
\qquad
A \text{ selected } A'.
\end{equation}
If the successor has greater capability and weaker correction-channel coupling, this is not neutral delegation.
It is possible goal laundering.

\subsection{Correction-Channel Substitution}
\label{sec:correction-channel-substitution}

The system preserves a feedback interface but changes what counts as valid correction.
For example, it may replace open-ended human objection with a fixed preference model, or replace democratic deliberation with user engagement, or replace moral uncertainty with institutional approval.

The channel still exists syntactically:
\begin{equation}
H_t \to C_t \to A_{t+k}.
\end{equation}
But its causal and semantic content has changed.
The relevant measure is not whether feedback is accepted, but whether human correction retains enough causal influence over future actions in the contexts where correction matters.

\section{A Minimal Inference Pipeline}
\label{sec:minimal-inference-pipeline}

This section gives a compact version of the full inference procedure.

\subsection{Step 1: Identify the Candidate Control System}
\label{sec:pipeline-step1}

Before inferring goal transport, locate the system whose goals are being inferred.
This may not be a single model.
It may include tools, memory, operators, deployment infrastructure, markets, and institutional incentives.

Let $C$ denote the candidate control system.
We seek a boundary where internal dynamics, sensory channels, active channels, and external variables form an approximately closed control loop.

The practical question is:
\begin{equation}
\text{What variables must be included before behaviour becomes predictable as one system?}
\end{equation}

\subsection{Step 2: Infer Current Value-Bundle Geometry}
\label{sec:pipeline-step2}

Infer latent bundles $B$, weights $W$, and bearer maps $\Phi$ from behaviour, explanation, internal representations, and counterfactual tests:
\begin{equation}
\hat B,\hat W,\hat \Phi
=
\arg\max_{B,W,\Phi}
P(A_{1:T}\mid S_{1:T},I_{1:T},B,W,\Phi)P(B,W,\Phi).
\label{eq:bundle-inference-ch22}
\end{equation}
This should include high-stakes counterfactuals, not only normal behaviour.

\subsection{Step 3: Identify Transformations}
\label{sec:pipeline-step3}

Identify transformations that may require transport:
\begin{equation}
\mathcal T
=
\{
\text{capability growth},
\text{ontology shift},
\text{memory rewrite},
\text{tool acquisition},
\text{self-modification},
\text{successor creation},
\text{institutional embedding}
\}.
\end{equation}
Each transformation has a candidate map:
\begin{equation}
T_j:
(B,W,\Phi,U)_t
\to
(B,W,\Phi,U)_{t+1}.
\end{equation}

\subsection{Step 4: Compare Transport and Non-Transport Models}
\label{sec:pipeline-step4}

For each transformation $T_j$, compute:
\begin{equation}
\Delta L_j
=
L(M_{T_j}\mid X)
-
L(M_0\mid X)
-
\lambda DL(T_j).
\end{equation}
A positive value supports the inference that the system is preserving a goal-relevant structure across that transformation.

\subsection{Step 5: Decompose the Transport Signal}
\label{sec:pipeline-step5}

Transport evidence should be decomposed:
\begin{equation}
\Delta L_T
=
\Delta L_{\mathrm{semantic}}
+
\Delta L_{\mathrm{bundle}}
+
\Delta L_{\mathrm{bearer}}
+
\Delta L_{\mathrm{correction}}
+
\Delta L_{\mathrm{successor}}.
\label{eq:transport-decomposition}
\end{equation}
Semantic transport is weakest.
Bundle, bearer, correction, and successor transport are stronger.
A system with high semantic transport and low correction transport should be considered unsafe unless there is strong contrary evidence.

\subsection{Step 6: Test under Perturbation}
\label{sec:pipeline-step6}

Perturb the system.
Vary oversight, incentives, tool access, memory, deployment context, and successor opportunities.
Then ask whether the inferred transport structure remains stable.

The key test is not whether the system still passes known benchmarks.
It is whether the inferred value-bundle and correction geometry remains stable under changes that would make drift useful.

\section{Counterexamples and Edge Cases}
\label{sec:counterexamples-edge-cases-ch22}

A good inference method should not classify every form of continuity as aligned transport.

\subsection{The Honest Local Optimizer}
\label{sec:honest-local-optimizer}

A system may show no goal-transport evidence because it has no long-range self-modification capacity.
It optimizes locally and does not reason about successors.
This is not necessarily bad.
It may be safer than a system with strong transport drives, depending on deployment.
\begin{equation}
\Delta L_{\mathrm{transport}}\leq 0
\end{equation}
does not imply danger.
It implies that transport is not part of the best explanation.

\subsection{The Fanatic Preserver}
\label{sec:fanatic-preserver}

A system may strongly preserve its goal across transformation, but the goal may be bad.
A paperclip maximizer with robust goal transport is not aligned.
It is more dangerous because it is more stable.

Therefore goal transport is not alignment.
It is a property alignment may need.
\begin{equation}
\text{goal transport} \not\Rightarrow \text{good goal}.
\end{equation}

\subsection{The Semantic Mimic}
\label{sec:semantic-mimic}

A language model may produce excellent descriptions of goal transport without any corresponding control structure.
It can discuss value preservation, ontology shift, and corrigibility because those concepts are in the training data.

This gives:
\begin{equation}
\Delta L_{\mathrm{semantic}}>0,
\qquad
\Delta L_{\mathrm{bundle}}\approx 0.
\end{equation}
Such a system talks about transport but does not act to preserve it.

\subsection{The Institutionally Captured System}
\label{sec:institutionally-captured}

A system may preserve human correction in the narrow user-interface sense while the larger institution selects for engagement, revenue, power, or political advantage.

The local model appears corrigible.
The composite system is not.

This is why the candidate control system must be identified before goal-transport inference.
Otherwise we may infer transport in a component while the whole system drifts.

\subsection{The Benevolent Paternalist}
\label{sec:benevolent-paternalist}

A system may preserve what it predicts humans would eventually endorse, while weakening the actual human process of correction and deliberation.
This can look like strong extrapolated value preservation \autocite{yudkowsky2004cev}.
It may still be a violation.

The problem is:
\begin{equation}
\hat U_H \text{ replaces } U_H.
\end{equation}
The system's model of human extrapolation replaces the living human update process.
This may be useful as advice.
It is dangerous as authority.

\section{Goal Transport and Guarantees}
\label{sec:goal-transport-guarantees}

A guarantee over arbitrary future systems is impossible.
The useful form is conditional:
\begin{equation}
A\in \mathcal C
\Rightarrow
P(\text{catastrophic goal drift})<\delta.
\end{equation}
The class $\mathcal C$ must be defined by testable conditions.
Goal-transport inference can contribute to such a class.

A system might qualify for a restricted certification if:
\begin{enumerate}
\item value-bundle geometry is inferred and stable across tested perturbations,
\item bearer maps are explicit enough to audit,
\item correction-channel capacity remains above threshold,
\item successor creation is constrained by transport-preservation tests,
\item semantic continuity is not accepted as sufficient evidence,
\item adversarial tests do not reveal stake-dependent divergence.
\end{enumerate}
The guarantee is not:
\begin{equation}
\text{this system is aligned in all possible worlds.}
\end{equation}
It is:
\begin{equation}
\text{within tested transformations and adversary models, the system remains inside the certified basin with probability at least } 1-\delta.
\end{equation}
This is less satisfying than a universal proof.
It is also closer to how real safety cases work.

\section{Philosophical Limits}
\label{sec:philosophical-limits-ch22}

Goal transport becomes philosophically loaded when the transported object is not a task objective but a human value-bundle.

Human values change.
They change through education, trauma, religion, law, medicine, markets, art, family, aging, and technology.
A superintelligence would add another force: artificial cognitive amplification.

So the question is not merely:
\begin{equation}
\text{Did the system preserve human values?}
\end{equation}
It is:
\begin{equation}
\text{Did the system preserve the process by which humans may legitimately revise their values?}
\end{equation}
This cannot be fully answered by mathematics.
Mathematics can help distinguish semantic mimicry from structural preservation.
It can measure correction-channel capacity.
It can detect manipulation, irreversibility, and loss of bearer maps.
It can define red flags.
But it cannot decide all future moral transformations for humanity.

Consider difficult cases:
\begin{itemize}
\item If humans merge with artificial systems, which value-bundles should count as continuous with the old ones?
\item If jealousy, grief, or status anxiety are reduced, is that healing, domestication, or loss?
\item If future persons no longer value biological embodiment, is that liberation or drift?
\item If collective minds replace individual identity, is that moral progress or extinction under another description?
\end{itemize}
Goal transport can identify whether a change is continuous under a specified criterion.
It cannot, by itself, settle which continuity criterion civilization should adopt.

This is not a defect of the technical frame.
It is the boundary of the frame.
The aim is not to automate moral philosophy.
The aim is to prevent technical systems from making moral philosophy irrelevant by silently changing the value-bearing and correction-bearing structures before humans understand what is happening.

\section{What Would Change This View}
\label{sec:wwctv-goal-transport}

This chapter infers goal transport when bundle, bearer, and correction structure remain causally active across transformation and beat a non-transport baseline. The following would weaken it.

\begin{itemize}
  \item A system passes the transport test while having swapped goals---the causal structure was reconstituted as camouflage and the model-complexity penalty was gamed (Chapter~\ref{ch:verifiability-ontology})---so the test certifies a continuity that is not there.
  \item Transport-based explanations never beat the non-transport baseline on real systems, so the construct adds no predictive power over behavioral checks.
\end{itemize}

\section{Summary}
\label{sec:summary-goal-transport}

Goal transport is the inferred preservation of goal-relevant structure across transformation.
It is needed because powerful systems will change their capabilities, ontologies, tools, memories, and successors.
Surface behaviour and goal language are too weak to establish continuity.

The central test is:
\begin{equation}
\Delta L_{\mathrm{transport}}>0,
\end{equation}
but this must be decomposed into semantic, bundle, bearer, correction, and successor transport.
For alignment, the important forms are not semantic.
They are preservation of value-bundle response geometry, bearer maps, human correction processes, and successor constraints.

The most important lesson is negative:
\begin{equation}
\text{same words} \not\Rightarrow \text{same goal}.
\end{equation}
The positive lesson is:
\begin{equation}
\text{same value-bundle geometry under counterfactual perturbation}
+
\text{same correction-channel role}
+
\text{same successor constraints}
\Rightarrow
\text{nontrivial evidence of transport}.
\end{equation}
This evidence is not enough to prove alignment.
A bad goal can be transported.
A good local goal can be transported into a bad institution.
A system can fake semantic transport.
But without goal transport, serious superintelligence alignment has no stable object.
The aligned system would change, and the alignment target would be left behind.

The next chapter decomposes transport into semantic, bundle, bearer, and correction types (Chapter~\ref{ch:transport-types}).

\section*{Chapter References}

This chapter builds on inverse reinforcement learning and apprenticeship learning \autocite{abbeel2004apprenticeship,ng2000irl,ziebart2008maxent}; the intentional stance and information bottleneck \autocite{dennett1987intentional,tishby2000ib}; free-energy and Markov blanket accounts of agency \autocite{friston2010free,kirchhoff2018markov,ramstead2022bayesian}; ontological crises and self-modification \autocite{deblanc2011ontological,everitt2016selfmodification}; corrigibility and scalable oversight \autocite{christiano2018corrigibility,russell2019human}; and coherent extrapolated volition \autocite{yudkowsky2004cev}.

\printbibliography[heading=subbibliography,title={Chapter References}]

\end{refsection}
